\section{Préschémas réduits~; condition de séparation}
\label{section:I.5-fr}

\subsection{Préschémas réduits.}
\label{subsection:I.5.1-fr}

\begin{proposition}[5.1.1]
\label{I.5.1.1-fr}
Soient $(X,\mathscr{O}_X)$ un préschéma, $\mathscr{B}$ une
$\mathscr{O}_X$-Algèbre quasi-cohérente. Il existe un
$\mathscr{O}_X$-Module quasi-cohérent et un seul $\mathscr{N}$ dont la fibre
$\mathscr{N}_x$ en tout $x\in X$ soit le nilradical de l'anneau
$\mathscr{B}_x$. Lorsque $X$ est affine, et par suite
$\mathscr{B}=\widetilde{B}$, où $B$ est une algèbre sur $A(X)$, on a
$\mathscr{N}=\widetilde{\mathfrak{N}}$, où $\mathfrak{N}$ est le nilradical
de $B$.
\end{proposition}

\oldpage[I]{128}
La question étant locale, on est ramené à démontrer la dernière assertion. On
sait que $\widetilde{\mathfrak{N}}$ est un $\mathscr{O}_X$-Module
quasi-cohérent (\hyperref[I.1.4.1-fr]{1.4.1}) et que sa fibre au point
$x\in X$ est l'idéal $\mathfrak{N}_x$ de l'anneau de fractions $B_x$~; tout
revient à prouver que le nilradical de $B_x$ est contenu dans
$\mathfrak{N}_x$, l'inclusion opposée étant évidente. Or, soit $z/s$ un
élément du nilradical de $B_x$, avec $z\in B$, $s\notin\mathfrak{j}_x$~; par
hypothèse, il existe un entier $k$ tel que $(z/s)^k=0$, ce qui signifie qu'il
existe $t\notin\mathfrak{j}_x$ tel que $tz^k=0$. On en conclut $(tz)^k=0$, et
par suite $z/s=(tz)/(ts)$ appartient à $\mathfrak{N}_x$.

Nous dirons que le $\mathscr{O}_X$-Module quasi-cohérent $\mathscr{N}$ ainsi
défini est le \emph{Nilradical} de la $\mathscr{O}_X$-Algèbre $\mathscr{B}$~;
nous noterons en particulier $\mathscr{N}_X$ le nilradical de
$\mathscr{O}_X$.

\begin{corollary}[5.1.2]
\label{I.5.1.2-fr}
Soit $X$ un préschéma~; le sous-préschéma fermé de $X$ défini par le faisceau
d'idéaux $\mathscr{N}_X$ est le seul sous-préschéma réduit
(\hyperref[0.4.1.4-fr]{0, 4.1.4}) de $X$ ayant $X$ pour espace sous-jacent~;
c'est aussi le plus petit sous-préschéma de $X$ ayant $X$ pour espace
sous-jacent.
\end{corollary}

Comme le faisceau structural du sous-préschéma fermé défini $Y$ par
$\mathscr{N}_X$ est $\mathscr{O}_X/\mathscr{N}_X$, il est immédiat que $Y$
est réduit et a $X$ pour espace sous-jacent, puisque
$\mathscr{N}_x\ne\mathscr{O}_x$ pour tout $x\in X$. Pour démontrer les autres
assertions, remarquons qu'un sous-préschéma $Z$ de $X$ ayant $X$ pour espace
sous-jacent est défini par un faisceau d'idéaux $\mathscr{J}$
(\hyperref[I.4.1.3-fr]{4.1.3}) tel que
$\mathscr{J}_x\ne\mathscr{O}_x$ pour tout $x\in X$. On peut se borner au cas
où $X$ est affine, soit $X=\operatorname{Spec}(A)$ et
$\mathscr{J}=\widetilde{\mathfrak{J}}$, où $\mathfrak{J}$ est un idéal de
$A$~; alors, pour tout $x\in X$, on a
$\mathfrak{J}\subset\mathfrak{j}_x$, donc $\mathfrak{J}$ est contenu dans
tous les idéaux premiers de $A$, c'est-à-dire dans leur intersection
$\mathfrak{N}$, nilradical de $A$. Cela prouve que $Y$ est le plus petit
sous-préschéma de $X$ ayant $X$ comme espace sous-jacent
(\hyperref[I.4.1.9-fr]{4.1.9})~; en outre, si $Z$ est distinct de $Y$, on a
nécessairement $\mathscr{J}_x\ne\mathscr{N}_x$ pour un $x\in X$ au moins, et
par suite (\hyperref[I.5.1.1-fr]{5.1.1}) $Z$ n'est pas réduit.

\begin{definition}[5.1.3]
\label{I.5.1.3-fr}
On appelle préschéma réduit associé à un préschéma $X$ et on note
$X_{\mathrm{red}}$ l'unique sous-préschéma réduit de $X$ ayant $X$ pour
espace sous-jacent.
\end{definition}

Dire qu'un préschéma $X$ est réduit signifie donc que
$X=X_{\mathrm{red}}$.

\begin{proposition}[5.1.4]
\label{I.5.1.4-fr}
Pour que le spectre premier d'un anneau $A$ soit un préschéma réduit (resp.
intègre) (\hyperref[I.2.1.7-fr]{2.1.7}), il faut et il suffit que $A$ soit un
anneau réduit (resp. intègre).
\end{proposition}

En effet, il résulte aussitôt de (\hyperref[I.5.1.1-fr]{5.1.1}) que la
condition $\mathscr{N}=(0)$ est nécessaire et suffisante pour que
$X=\operatorname{Spec}(A)$ soit réduit~; l'assertion relative aux anneaux
intègres est alors une conséquence de
(\hyperref[I.1.1.13-fr]{1.1.13}).

Comme tout anneau de fractions $\ne\{0\}$ d'un anneau intègre est intègre, il
résulte de (\hyperref[I.5.1.4-fr]{5.1.4}) que pour tout préschéma
\emph{localement intègre} $X$, $\mathscr{O}_x$ est un anneau \emph{intègre}
pour tout $x\in X$. La réciproque est vraie lorsque l'espace sous-jacent à
$X$ est \emph{localement noethérien}~: en effet $X$ est alors réduit, et si
$U$ est un ouvert affine de $X$, qui soit un espace noethérien, $U$ n'a qu'un
nombre fini de composantes irréductibles, donc son anneau $A$ n'a qu'un
nombre fini d'idéaux premiers minimaux
(\hyperref[I.1.1.14-fr]{1.1.14}). Si deux de ces composantes $U_i$ avaient un
point commun $x$, $\mathscr{O}_x$ aurait au moins deux idéaux premiers
minimaux distincts, donc ne serait pas intègre~; les $U_i$ sont par suite des
ouverts deux à deux disjoints, et chacun d'eux est donc intègre.

\begin{env}[5.1.5]
\label{I.5.1.5-fr}
Soit $f=(\psi,\theta):X\to Y$ un morphisme de préschémas~; l'homomor-
\oldpage[I]{129}
phisme $\theta_x^\sharp:\mathscr{O}_{\psi(x)}\to\mathscr{O}_x$ applique tout
élément nilpotent de $\mathscr{O}_{\psi(x)}$ sur un élément nilpotent de
$\mathscr{O}_x$~; par passage aux quotients, on déduit donc de $\theta^\sharp$
un homomorphisme
\[
  \omega:\psi^*(\mathscr{O}_Y/\mathscr{N}_Y)
  \longrightarrow \mathscr{O}_X/\mathscr{N}_X
\]
il est clair que pour tout $x\in X$,
$\omega_x:\mathscr{O}_{\psi(x)}/\mathscr{N}_{\psi(x)}
\to\mathscr{O}_x/\mathscr{N}_x$ est un homomorphisme local, donc
$(\psi,\omega^b)$ est un morphisme de préschémas
$X_{\mathrm{red}}\to Y_{\mathrm{red}}$, que nous noterons $f_{\mathrm{red}}$
et appellerons le morphisme \emph{réduit} associé à $f$. Il est immédiat que
pour deux morphismes $f:X\to Y$, $g:Y\to Z$, on a
$(g\circ f)_{\mathrm{red}}=g_{\mathrm{red}}\circ f_{\mathrm{red}}$, donc on a
défini $X_{\mathrm{red}}$ comme \emph{foncteur covariant} en $X$.

La définition précédente montre que le diagramme
\[
  \xymatrix{
    X_{\mathrm{red}}\ar[r]^{f_{\mathrm{red}}}\ar[d] &
    Y_{\mathrm{red}}\ar[d]\\
    X\ar[r]^f & Y
  }
\]
est commutatif, les flèches verticales étant les morphismes d'injection~; en
d'autres termes, $X_{\mathrm{red}}\to X$ est un morphisme \emph{fonctoriel}.
On notera en particulier que si $X$ est \emph{réduit}, tout morphisme
$f:X\to Y$ se factorise en
$X\xrightarrow{f_{\mathrm{red}}}Y_{\mathrm{red}}\to Y$~; en d'autres termes,
$f$ est \emph{majoré} par le morphisme d'injection
$Y_{\mathrm{red}}\to Y$.
\end{env}

\begin{proposition}[5.1.6]
\label{I.5.1.6-fr}
Soit $f:X\to Y$ un morphisme~; si $f$ est surjectif (resp. radiciel, une
immersion, une immersion fermée, une immersion ouverte, une immersion locale,
un isomorphisme local), il en est de même de $f_{\mathrm{red}}$. Inversement,
si $f_{\mathrm{red}}$ est surjectif (resp. radiciel), il en est de même de
$f$.
\end{proposition}

La proposition est triviale si $f$ est surjectif~; si $f$ est radiciel, elle
résulte de ce que pour tout $x\in X$, le corps $k(x)$ est le même pour les
préschémas $X$ et $X_{\mathrm{red}}$
(\hyperref[I.3.5.8-fr]{3.5.8}). Enfin, si $f=(\psi,\theta)$ est une immersion,
une immersion fermée ou une immersion locale (resp. une immersion ouverte, ou
un isomorphisme local), la proposition résulte de ce que si
$\theta_x^\sharp$ est surjectif (resp. bijectif), il en est de même de
l'homomorphisme obtenu en passant aux quotients par les nilradicaux de
$\mathscr{O}_{\psi(x)}$ et $\mathscr{O}_x$
(\hyperref[I.5.1.2-fr]{5.1.2} et \hyperref[I.4.2.2-fr]{4.2.2})
(cf. (\hyperref[I.5.5.12-fr]{5.5.12})).

\begin{proposition}[5.1.7]
\label{I.5.1.7-fr}
Si $X$, $Y$ sont deux $S$-préschémas, les préschémas
$X_{\mathrm{red}}\times_{S_{\mathrm{red}}}Y_{\mathrm{red}}$ et
$X_{\mathrm{red}}\times_S Y_{\mathrm{red}}$ sont identiques, et
s'identifient canoniquement à un sous-préschéma de $X\times_S Y$ ayant même
espace sous-jacent que ce produit.
\end{proposition}

L'identification canonique de $X_{\mathrm{red}}\times_S Y_{\mathrm{red}}$ à
un sous-préschéma de $X\times_S Y$ ayant même espace sous-jacent résulte de
(\hyperref[I.4.3.1-fr]{4.3.1}). D'autre part, si $\varphi$ et $\psi$ sont les
morphismes structuraux $X_{\mathrm{red}}\to S$, $Y_{\mathrm{red}}\to S$, ils
se factorisent par $S_{\mathrm{red}}$
(\hyperref[I.5.1.5-fr]{5.1.5}), et comme $S_{\mathrm{red}}\to S$ est un
monomorphisme, la première assertion résulte de
(\hyperref[I.3.2.4-fr]{3.2.4}).

\begin{corollary}[5.1.8]
\label{I.5.1.8-fr}
Les préschémas $(X\times_S Y)_{\mathrm{red}}$ et
$(X_{\mathrm{red}}\times_{S_{\mathrm{red}}}Y_{\mathrm{red}})_{\mathrm{red}}$
s'identifient canoniquement.
\end{corollary}

Cela résulte de (\hyperref[I.5.1.2-fr]{5.1.2} et
\hyperref[I.5.1.7-fr]{5.1.7}).

On notera que si $X$ et $Y$ sont des préschémas réduits, il n'en est pas
nécessairement de même de $X\times_S Y$, car le produit tensoriel de deux
algèbres réduites peut avoir des éléments nilpotents.

\begin{proposition}[5.1.9]
\label{I.5.1.9-fr}
\oldpage[I]{130}
Soient $X$ un préschéma, $\mathscr{J}$ un faisceau quasi-cohérent d'idéaux de
$\mathscr{O}_X$ tel que $\mathscr{J}^n=0$ pour un entier $n>0$. Soit $X_0$ le
sous-préschéma fermé $(X,\mathscr{O}_X/\mathscr{J})$ de $X$~; pour que $X$
soit un schéma affine, il faut et il suffit que $X_0$ le soit.
\end{proposition}

La condition étant évidemment nécessaire, prouvons qu'elle est suffisante.
Si on pose $X_k=(X,\mathscr{O}_X/\mathscr{J}^{k+1})$, tout revient à prouver
par récurrence sur $k$ que les $X_k$ sont affines, donc on est ramené au cas
où $\mathscr{J}^2=0$. Posons
\[
  A=\Gamma(X,\mathscr{O}_X),\qquad
  A_0=\Gamma(X_0,\mathscr{O}_{X_0})
      =\Gamma(X,\mathscr{O}_X/\mathscr{J}).
\]
On déduit de l'homomorphisme canonique
$\mathscr{O}_X\to\mathscr{O}_X/\mathscr{J}$ un homomorphisme d'anneaux
$\varphi:A\to A_0$. Nous verrons ci-dessous que $\varphi$ est
\emph{surjectif}, de sorte que la suite
\[
  0\to\Gamma(X,\mathscr{J})\to\Gamma(X,\mathscr{O}_X)
  \to\Gamma(X,\mathscr{O}_X/\mathscr{J})\to 0
  \tag{5.1.9.1}\label{I.5.1.9.1-fr}
\]
est \emph{exacte}. Supposons ce point établi, et montrons que cela entraîne la
proposition. Notons que $\mathfrak{K}=\Gamma(X,\mathscr{J})$ est un idéal de
carré nul dans $A$, et est donc un module sur $A_0=A/\mathfrak{K}$. Par
hypothèse, on a $X_0=\operatorname{Spec}(A_0)$, et comme les espaces
topologiques sous-jacents $X_0$ et $X$ sont identiques,
$\mathfrak{K}=\Gamma(X_0,\mathscr{J})$~; par ailleurs, comme
$\mathscr{J}^2=0$, $\mathscr{J}$ est un
$(\mathscr{O}_X/\mathscr{J})$-Module quasi-cohérent, donc on a
$\mathscr{J}\simeq\widetilde{\mathfrak{K}}$ et
$\mathfrak{K}_x=\mathscr{J}_x$ pour tout $x\in X_0$
(\hyperref[I.1.4.1-fr]{1.4.1}).

Cela étant, soit $X'=\operatorname{Spec}(A)$, et considérons le morphisme
$f=(\psi,\theta):X\to X'$ de préschémas correspondant à l'application
identique $A\to\Gamma(X,\mathscr{O}_X)$
(\hyperref[I.2.2.4-fr]{2.2.4}). Pour tout ouvert affine $V$ dans $X$, le
diagramme
\[
  \xymatrix{
    A\ar[r]\ar[d] & \Gamma(V,\mathscr{O}_X|V)\ar[d]\\
    A_0=A/\mathfrak{K}\ar[r] & \Gamma(V,\mathscr{O}_{X_0}|V)
  }
\]
est commutatif, d'où on conclut que le diagramme
\[
  \xymatrix{
    X' & X\ar[l]_f\\
    X'_0\ar[u]^{j'} & X_0\ar[l]_{f_0}\ar[u]_j
  }
\]
est commutatif, $X'_0$ étant le sous-préschéma fermé de $X'$ défini par le
faisceau quasi-cohérent d'idéaux $\widetilde{\mathfrak{K}}$, $j$, $j'$ les
morphismes d'injection canoniques. Mais puisque $X_0$ est affine, $f_0$ est un
isomorphisme et comme les applications continues sous-jacentes à $j$ et $j'$
sont les applications identiques, on voit tout d'abord que
$\psi:X\to X'$ est un homéomorphisme. En outre, la relation
$\mathfrak{K}_x=\mathscr{J}_x$ montre que la restriction de
$\theta^\sharp:\psi^*(\mathscr{O}_{X'})\to\mathscr{O}_X$ est un
\emph{isomorphisme} de $\psi^*(\widetilde{\mathfrak{K}})$ sur $\mathscr{J}$~;
d'autre part, par passage aux quotients, $\theta^\sharp$ donne un
\emph{isomorphisme}
$\psi^*(\mathscr{O}_{X'}/\widetilde{\mathfrak{K}})
\to\mathscr{O}_X/\mathscr{J}$, puisque $f_0$ est un isomorphisme~; on en
conclut aussitôt par le lemme des 5 (M, I, 1.1) que $\theta^\sharp$ est
lui-même un isomorphisme, donc que $f$ est un \emph{isomorphisme}, et par suite
que $X$ est affine.

Tout revient donc à démontrer l'exactitude de
(\hyperref[I.5.1.9.1-fr]{5.1.9.1}), ce qui résultera de

\oldpage[I]{131}
$H^1(X,\mathscr{J})=0$. Or,
$H^1(X,\mathscr{J})=H^1(X_0,\mathscr{J})$ et nous avons vu que
$\mathscr{J}$ est un $\mathscr{O}_{X_0}$-Module quasi-cohérent. Notre
assertion résultera donc du

\begin{lemma}[5.1.9.2]
\label{I.5.1.9.2-fr}
Si $Y$ est un schéma affine et $\mathscr{F}$ un $\mathscr{O}_Y$-Module
quasi-cohérent, on a $H^1(Y,\mathscr{F})=0$.
\end{lemma}

Ce lemme sera démontré au chap. III, § 1, comme conséquence du théorème plus
général que $H^i(Y,\mathscr{F})=0$ pour tout $i>0$. Pour en donner une
démonstration indépendante, observons que $H^1(Y,\mathscr{F})$ s'identifie au
module $\operatorname{Ext}^1_{\mathscr{O}_Y}
(Y;\mathscr{O}_Y,\mathscr{F})$ des classes d'extensions du
$\mathscr{O}_Y$-Module $\mathscr{O}_Y$ par le $\mathscr{O}_Y$-Module
$\mathscr{F}$ (T, 4.2.3)~; tout revient donc à prouver qu'une telle extension
$\mathscr{G}$ est triviale. Or, pour tout $y\in Y$, il y a un voisinage $V$
de $y$ dans $Y$ tel que $\mathscr{G}|V$ soit isomorphe à
$\mathscr{F}|Y\oplus\mathscr{O}_Y|V$
(\hyperref[0.5.4.9-fr]{0, 5.4.9})~; on en conclut que $\mathscr{G}$ est un
$\mathscr{O}_Y$-Module quasi-cohérent. Si $A$ est l'anneau de $Y$, on a donc
$\mathscr{F}=\widetilde{M}$, $\mathscr{G}=\widetilde{N}$, où $M$ et $N$ sont
des $A$-modules, et par hypothèse $N$ est une extension du $A$-module $A$ par
le $A$-module $M$ (\hyperref[I.1.3.11-fr]{1.3.11}). Comme cette extension est
nécessairement triviale, le lemme est démontré, et par suite aussi
(\hyperref[I.5.1.9-fr]{5.1.9}).

\begin{corollary}[5.1.10]
\label{I.5.1.10-fr}
Soit $X$ un préschéma tel que $\mathscr{N}_X$ soit nilpotent. Pour que $X$
soit un schéma affine, il faut et il suffit que $X_{\mathrm{red}}$ le soit.
\end{corollary}

\subsection{Existence d'un sous-préschéma d'espace sous-jacent donné.}
\label{subsection:I.5.2-fr}

\begin{proposition}[5.2.1]
\label{I.5.2.1-fr}
Pour tout sous-espace localement fermé $Y$ de l'espace sous-jacent à un
préschéma $X$, il existe un sous-préschéma réduit et un seul de $X$ ayant $Y$
pour espace sous-jacent.
\end{proposition}

L'unicité résultant de (\hyperref[I.5.1.2-fr]{5.1.2}), il reste à prouver
l'existence du sous-préschéma en question.

Si $X$ est affine d'anneau $A$, et $Y$ fermé dans $X$, la proposition est
immédiate~: $\mathfrak{j}(Y)$ est le plus grand idéal
$\mathfrak{a}\subset A$ tel que $V(\mathfrak{a})=Y$, et il est égal à sa
racine (\hyperref[I.1.1.4-fr]{1.1.4 (i)}), donc
$A/\mathfrak{j}(Y)$ est un anneau réduit.

Dans le cas général, pour tout ouvert affine $U\subset X$ tel que $U\cap Y$
soit fermé dans $U$, considérons le sous-préschéma fermé $Y_U$ de $U$ défini
par le faisceau d'idéaux associé à l'idéal $\mathfrak{j}(U\cap Y)$ de $A(U)$,
qui est réduit. Montrons que, si $V$ est un ouvert affine de $X$ contenu dans
$U$, $Y_V$ est \emph{induit} par $Y_U$ sur $V\cap Y$~; or, ce préschéma
induit est un sous-préschéma fermé de $V$ qui est réduit et a $V\cap Y$ comme
espace sous-jacent~; l'unicité de $Y_V$ entraîne donc notre assertion.

\begin{proposition}[5.2.2]
\label{I.5.2.2-fr}
Soient $X$ un préschéma réduit, $f:X\to Y$ un morphisme, $Z$ un
sous-préschéma fermé de $Y$ tel que $f(X)\subset Z$~; alors $f$ se factorise
en $X\xrightarrow{g}Z\xrightarrow{j}Y$, où $j$ est le morphisme d'injection.
\end{proposition}

Il résulte de l'hypothèse que le sous-préschéma fermé $f^{-1}(Z)$ de $X$ a
pour espace sous-jacent $X$ tout entier
(\hyperref[I.4.4.1-fr]{4.4.1})~; comme $X$ est réduit, ce sous-préschéma fermé
coïncide avec $X$ (\hyperref[I.5.1.2-fr]{5.1.2}), et la proposition résulte
donc de (\hyperref[I.4.4.1-fr]{4.4.1}).

\begin{corollary}[5.2.3]
\label{I.5.2.3-fr}
Soit $X$ un sous-préschéma réduit d'un préschéma $Y$~; si $Z$ est le
sous-préschéma fermé réduit de $Y$ ayant pour espace sous-jacent
$\overline{X}$, $X$ est un sous-préschéma induit sur un ouvert de $Z$.
\end{corollary}

\oldpage[I]{132}
Il y a en effet un ouvert $U$ de $Y$ tel que
$X=U\cap\overline{X}$~; comme $X$ est un sous-préschéma réduit de $Z$ en
vertu de (\hyperref[I.5.2.2-fr]{5.2.2}), le sous-préschéma $X$ est induit par
$Z$ sur le sous-espace ouvert $X$ en vertu de l'unicité
(\hyperref[I.5.2.1-fr]{5.2.1}).

\begin{corollary}[5.2.4]
\label{I.5.2.4-fr}
Soient $f:X\to Y$ un morphisme, $X'$ (resp. $Y'$) un sous-préschéma fermé de
$X$ (resp. $Y$) défini par un faisceau d'idéaux quasi-cohérent $\mathscr{J}$
(resp. $\mathscr{K}$) de $\mathscr{O}_X$ (resp. $\mathscr{O}_Y$). Supposons
que $X'$ soit réduit et que $f(X')\subset Y'$. Alors on a
$f^*(\mathscr{K})\mathscr{O}_X\subset\mathscr{J}$.
\end{corollary}

Comme la restriction de $f$ à $X'$ se factorise en
$X'\to Y'\to Y$ d'après (\hyperref[I.5.2.2-fr]{5.2.2}), il suffit d'appliquer
(\hyperref[I.4.4.6-fr]{4.4.6}).

\subsection{Diagonale; graphe d'un morphisme.}
\label{subsection:I.5.3-fr}

\begin{env}[5.3.1]
\label{I.5.3.1-fr}
Soit $X$ un $S$-préschéma~; on appelle \emph{morphisme diagonal} de $X$ dans
$X\times_S X$, et on note $\Delta_{X|S}$, ou $\Delta_X$, ou même $\Delta$ si
aucune confusion n'est possible, le $S$-morphisme $(1_X,1_X)_S$, autrement
dit, l'unique $S$-morphisme $\Delta_X$ tel que
\[
  p_1\circ\Delta_X=p_2\circ\Delta_X=1_X
  \tag{5.3.1.1}\label{I.5.3.1.1-fr}
\]
en désignant par $p_1$, $p_2$ les projections de $X\times_S X$
(déf. \hyperref[I.3.2.1-fr]{3.2.1}). Si $f:T\to X$, $g:T\to Y$ sont deux
$S$-morphismes, on vérifie aussitôt que
\[
  (f,g)_S=(f\times_S g)\circ\Delta_{T|S}
  \tag{5.3.1.2}\label{I.5.3.1.2-fr}
\]

Le lecteur observera que la définition précédente et les résultats énoncés
dans les n\textsuperscript{os} (\hyperref[I.5.3.1-fr]{5.3.1}) à (5.3.8) sont
valables dans toute catégorie, \emph{pourvu que les produits qui y figurent
existent dans cette catégorie}.
\end{env}

\begin{proposition}[5.3.2]
\label{I.5.3.2-fr}
Soient $X$, $Y$ deux $S$-préschémas~; si on identifie canoniquement le produit
$(X\times Y)\times(X\times Y)$ à $(X\times X)\times(Y\times Y)$, le morphisme
$\Delta_{X\times Y}$ s'identifie à $\Delta_X\times\Delta_Y$.
\end{proposition}

En effet, si $p_1$, $q_1$ sont les premières projections
$X\times X\to X$, $Y\times Y\to Y$, la première projection
$(X\times Y)\times(X\times Y)\to X\times Y$ s'identifie à $p_1\times q_1$,
et l'on a
\[
  (p_1\times q_1)\circ(\Delta_X\times\Delta_Y)
  =(p_1\circ\Delta_X)\times(q_1\circ\Delta_Y)=1_{X\times Y}
\]
même raisonnement pour les secondes projections.

\begin{corollary}[5.3.4]
\label{I.5.3.4-fr}
Pour toute extension $S'\to S$ du préschéma de base,
$\Delta_{X_{(S')}}$ s'identifie canoniquement à $(\Delta_X)_{(S')}$.
\end{corollary}

Il suffit de remarquer que $(X\times_S X)_{(S')}$ s'identifie canoniquement à
$X_{(S')}\times_{S'}X_{(S')}$
(\hyperref[I.3.3.10-fr]{3.3.10}).

\begin{proposition}[5.3.5]
\label{I.5.3.5-fr}
Soient $X$, $Y$ deux $S$-préschémas, $\varphi:S\to T$ un morphisme, faisant de
tout $S$-préschéma un $T$-préschéma. Soient $f:X\to S$, $Y\to S$ les
morphismes structuraux, $p$, $q$ les projections de $X\times_S Y$,
$\pi=f\circ p=g\circ q$ le morphisme structural $X\times_S Y\to S$. Alors, le
diagramme
\[
  \xymatrix{
    X\times_S Y\ar[r]^{(p,q)_T}\ar[d]_\pi &
    X\times_T Y\ar[d]^{f\times_T g}\\
    S\ar[r]^{\Delta_{S|T}} &
    S\times_T S
  }
  \tag{5.3.5.1}\label{I.5.3.5.1-fr}
\]
\oldpage[I]{133}
est commutatif, et identifie $X\times_S Y$ au produit des
$(S\times_T S)$-préschémas $S$ et $X\times_T Y$, les projections
s'identifiant à $\pi$ et $(p,q)_T$.
\end{proposition}

En vertu de (\hyperref[I.3.4.3-fr]{3.4.3}), on est ramené à démontrer la
proposition correspondante \emph{dans la catégorie des ensembles}, en
remplaçant $X$, $Y$, $S$ par $X(Z)_T$, $Y(Z)_T$, $S(Z)_T$, $Z$ étant un
$T$-préschéma arbitraire. Mais pour la catégorie des ensembles, la vérification
est immédiate et laissée au lecteur.

\begin{corollary}[5.3.6]
\label{I.5.3.6-fr}
Le morphisme $(p,q)_T$ s'identifie (en posant $P=S\times_T S$) à
$1_{X\times_T Y}\times_P\Delta_S$.
\end{corollary}

Cela résulte de (\hyperref[I.5.3.5-fr]{5.3.5}) et
(\hyperref[I.3.3.4-fr]{3.3.4}).

\begin{corollary}[5.3.7]
\label{I.5.3.7-fr}
Si $f:X\to Y$ est un $S$-morphisme, le diagramme
\[
  \xymatrix{
    X\ar[r]^{(1_X,f)_S}\ar[d]_f &
    X\times_S Y\ar[d]^{f\times_S 1_Y}\\
    Y\ar[r]^{\Delta_Y} &
    Y\times_S Y
  }
\]
est commutatif, et identifie $X$ au produit des
$(Y\times_S Y)$-préschémas $Y$ et $X\times_S Y$.
\end{corollary}

Il suffit d'appliquer (\hyperref[I.5.3.5-fr]{5.3.5}) en remplaçant $S$ par
$Y$ et $T$ par $S$, et remarquant que $X\times_Y Y=X$
(\hyperref[I.3.3.3-fr]{3.3.3}).

\begin{proposition}[5.3.8]
\label{I.5.3.8-fr}
Pour que $f:X\to Y$ soit un monomorphisme de préschémas, il faut et il suffit
que $\Delta_{X|Y}$ soit un isomorphisme de $X$ sur $X\times_Y X$.
\end{proposition}

En effet, dire que $f$ est un monomorphisme signifie que pour tout
$Y$-préschéma $Z$, l'application correspondante
$f':X(Z)_Y\to Y(Z)_Y$ est une injection, et comme $Y(Z)_Y$ est réduit à un
élément, cela signifie qu'il en est de même de $X(Z)_Y$. Mais cela s'exprime
aussi en disant que $X(Z)_Y\times X(Z)_Y$ est canoniquement isomorphe à
$X(Z)_Y$, et le premier de ces ensembles étant $(X\times_Y X)(Z)_Y$
(\hyperref[I.3.4.3.1-fr]{3.4.3.1}), cela signifie que $\Delta_{X|Y}$ est un
isomorphisme.

\begin{proposition}[5.3.9]
\label{I.5.3.9-fr}
Le morphisme diagonal $\Delta_X$ est une immersion de $X$ dans
$X\times_S X$.
\end{proposition}

En effet, comme les applications continues $p_1$ et $\Delta_X$ des espaces
sous-jacents sont telles que $p_1\circ\Delta_X$ soit l'identité, $\Delta_X$
est un homéomorphisme de $X$ sur $\Delta_X(X)$. De même, l'homomorphisme
composé
$\mathscr{O}_x\to\mathscr{O}_{\Delta_X(x)}\to\mathscr{O}_x$ des
homomorphismes correspondant à $p_1$ et à $\Delta_X$ étant l'identité,
l'homomorphisme correspondant à $\Delta_X$ est surjectif~; la proposition
résulte donc de (\hyperref[I.4.2.2-fr]{4.2.2}).

On dit que le sous-préschéma de $X\times_S X$ associé à l'immersion
$\Delta_X$ (\hyperref[I.4.2.1-fr]{4.2.1}) est \emph{la diagonale} de
$X\times_S X$.

\begin{corollary}[5.3.10]
\label{I.5.3.10-fr}
Sous les hypothèses de (\hyperref[I.5.3.5-fr]{5.3.5}), $(p,q)_T$ est une
immersion.
\end{corollary}

Cela résulte de (\hyperref[I.5.3.6-fr]{5.3.6}) et de
(\hyperref[I.4.3.1-fr]{4.3.1}).

On dit (sous les hypothèses de (\hyperref[I.5.3.5-fr]{5.3.5})) que $(p,q)_T$
est l'\emph{immersion canonique} de $X\times_S Y$ dans $X\times_T Y$.

\begin{corollary}[5.3.11]
\label{I.5.3.11-fr}
Soient $X$, $Y$ deux $S$-préschémas, $f:X\to Y$ un $S$-morphisme~; alors le
morphisme graphe $\Gamma_f=(1_X,f)_S$ de $f$
(\hyperref[I.3.3.14-fr]{3.3.14}) est une immersion de $X$ dans
$X\times_S Y$.
\end{corollary}

C'est le cas particulier du cor. (\hyperref[I.5.3.10-fr]{5.3.10}) où l'on
remplace $S$ par $Y$ et $T$ par $S$
(cf. (\hyperref[I.5.3.7-fr]{5.3.7})).

\oldpage[I]{134}
Le sous-préschéma de $X\times_S Y$ associé à l'immersion $\Gamma_f$
(\hyperref[I.4.2.1-fr]{4.2.1}) est appelé le \emph{graphe} du morphisme $f$~;
les sous-préschémas de $X\times_S Y$ qui sont des graphes de morphismes
$X\to Y$ se caractérisent par le fait que la restriction à un tel
sous-préschéma $G$ de la projection $p_1:X\times_S Y\to X$ est un
\emph{isomorphisme} $g$ de $G$ sur $X$~: $G$ est alors le graphe du morphisme
$p_2\circ g^{-1}$, où $p_2$ est la projection $X\times_S Y\to Y$.

Lorsqu'on prend en particulier $X=S$, les $S$-morphismes $S\to Y$, qui ne
sont autres que les $S$-sections de $Y$ (\hyperref[I.2.5.5-fr]{2.5.5}) sont
égaux à leurs morphismes graphes~; les sous-préschémas de $Y$ qui sont des
graphes de $S$-sections (autrement dit, ceux qui sont isomorphes à $S$ par la
restriction du morphisme structural $Y\to S$) s'appellent encore les
\emph{images} de ces sections, ou, par abus de langage, les
\emph{$S$-sections} de $Y$.

\begin{corollary}[5.3.12]
\label{I.5.3.12-fr}
Les hypothèses et notations étant celles de
(\hyperref[I.5.3.11-fr]{5.3.11}), pour tout morphisme $g:S'\to S$, soit $f'$
l'image réciproque de $f$ par $g$ (\hyperref[I.3.3.7-fr]{3.3.7})~; alors
$\Gamma_{f'}$ est l'image réciproque de $\Gamma_f$ par $g$.
\end{corollary}

C'est un cas particulier de la formule
(\hyperref[I.3.3.10.1-fr]{3.3.10.1}).

\begin{corollary}[5.3.13]
\label{I.5.3.13-fr}
Soient $f:X\to Y$, $g:Y\to Z$ deux morphismes~; si $g\circ f$ est une
immersion (resp. une immersion locale), il en est de même de $f$.
\end{corollary}

En effet, $f$ se factorise en
\[
  X\xrightarrow{\Gamma_f}X\times_Z Y\xrightarrow{p_2}Y.
\]
D'autre part, $p_2$ s'identifie à $(g\circ f)\times_Z 1_Y$
(\hyperref[I.3.3.4-fr]{3.3.4})~; si $g\circ f$ est une immersion (resp. une
immersion locale), il en est de même de $p_2$
(\hyperref[I.4.3.1-fr]{4.3.1} et \hyperref[I.4.5.5-fr]{4.5.5}), et comme
$\Gamma_f$ est une immersion (\hyperref[I.5.3.11-fr]{5.3.11}), on conclut
par (\hyperref[I.4.2.4-fr]{4.2.4}) (resp.
(\hyperref[I.4.5.5-fr]{4.5.5})).

\begin{corollary}[5.3.14]
\label{I.5.3.14-fr}
Soient $j:X\to Y$, $g:X\to Z$ deux $S$-morphismes. Si $j$ est une immersion
(resp. une immersion locale), il en est de même de $(j,g)_S$.
\end{corollary}

En effet, si $p:Y\times_S Z\to Y$ est la première projection, on a
$j=p\circ(j,g)_S$, et il suffit d'appliquer
(\hyperref[I.5.3.13-fr]{5.3.13}).

\begin{proposition}[5.3.15]
\label{I.5.3.15-fr}
Si $f:X\to Y$ est un $S$-morphisme, le diagramme
\[
  \xymatrix{
    X\ar[r]^{\Delta_X}\ar[d]_f &
    X\times_S X\ar[d]^{f\times_S f}\\
    Y\ar[r]^{\Delta_Y} &
    Y\times_S Y
  }
  \tag{5.3.15.1}\label{I.5.3.15.1-fr}
\]
est commutatif (en d'autres termes $\Delta_X$ est un morphisme fonctoriel
dans la catégorie des préschémas).
\end{proposition}

La vérification est immédiate et laissée au lecteur.

\begin{corollary}[5.3.16]
\label{I.5.3.16-fr}
Si $X$ est un sous-préschéma de $Y$, la diagonale $\Delta_X(X)$ s'identifie à
un sous-préschéma de $\Delta_Y(Y)$, dont l'espace sous-jacent s'identifie à
\[
  \Delta_Y(Y)\cap p_1^{-1}(X)=\Delta_Y(Y)\cap p_2^{-1}(X)
\]
($p_1$, $p_2$ projections de $Y\times_S Y$).
\end{corollary}

Appliquons (\hyperref[I.5.3.15-fr]{5.3.15}) au morphisme d'injection
$f:X\to Y$~; on sait alors que $f\times_S f$ est une immersion, identifiant
l'espace sous-jacent à $X\times_S X$ au sous-espace
$p_1^{-1}(X)\cap p_2^{-1}(X)$ de $Y\times_S Y$
(\hyperref[I.4.3.1-fr]{4.3.1})~; en outre, si
$z\in\Delta_Y(Y)\cap p_1^{-1}(X)$, on a $z=\Delta_Y(y)$

\oldpage[I]{135}
et $y=p_1(z)\in X$, donc $y=f(y)$, et $z=\Delta_Y(f(y))$ appartient à
$\Delta_X(X)$ en vertu de la commutativité du diagramme
(\hyperref[I.5.3.15.1-fr]{5.3.15.1}).

\begin{corollary}[5.3.17]
\label{I.5.3.17-fr}
Soient $f_1:Y\to X$, $f_2:Y\to X$ deux $S$-morphismes, $y$ un point de $Y$
tel que $f_1(y)=f_2(y)=x$ et que les homomorphismes $k(x)\to k(y)$
correspondant à $f_1$ et $f_2$ soient identiques. Alors, si
$f=(f_1,f_2)_S$, le point $f(y)$ appartient à la diagonale
$\Delta_{X|S}(X)$.
\end{corollary}

Les deux homomorphismes $k(x)\to k(y)$ correspondant à $f_i$ ($i=1,2$)
définissent deux $S$-morphismes
$g_i:\operatorname{Spec}(k(y))\to\operatorname{Spec}(k(x))$ tels que les
diagrammes
\[
  \xymatrix{
    \operatorname{Spec}(k(y))\ar[r]^{g_i}\ar[d] &
    \operatorname{Spec}(k(x))\ar[d]\\
    Y\ar[r]^{f_i} & X
  }
\]
soient commutatifs. Le diagramme
\[
  \xymatrix{
    \operatorname{Spec}(k(y))\ar[rr]^{(g_1,g_2)_S}\ar[d] & &
    \operatorname{Spec}(k(x))\times_S\operatorname{Spec}(k(x))\ar[d]\\
    Y\ar[rr]^{(f_1,f_2)_S} & & X\times_S X
  }
\]
est donc aussi commutatif. Or il résulte de l'égalité $g_1=g_2$ que l'image
par $(g_1,g_2)_S$ de l'unique point de $\operatorname{Spec}(k(y))$
appartient à la diagonale de
$\operatorname{Spec}(k(x))\times_S\operatorname{Spec}(k(x))$~; la conclusion
résulte donc de (\hyperref[I.5.3.15-fr]{5.3.15}).

\subsection{Morphismes et préschémas séparés.}
\label{subsection:I.5.4-fr}

\begin{definition}[5.4.1]
\label{I.5.4.1-fr}
On dit qu'un morphisme de préschémas $f:X\to Y$ est \emph{séparé} si le
morphisme diagonal $X\to X\times_Y X$ est une \emph{immersion fermée}~; on
dit alors aussi que $X$ est un \emph{préschéma séparé au-dessus de $Y$}, ou un
\emph{$Y$-schéma}. On dit qu'un préschéma $X$ est séparé s'il est séparé
au-dessus de $\operatorname{Spec}(\mathbf{Z})$~; on dit aussi alors que $X$
est un \emph{schéma} (cf. (\hyperref[I.5.5.7-fr]{5.5.7})).
\end{definition}

En vertu de (\hyperref[I.5.3.9-fr]{5.3.9}), pour que $X$ soit séparé au-dessus
de $Y$, il faut et il suffit que $\Delta_X(X)$ soit un \emph{sous-espace
fermé} de l'espace sous-jacent à $X\times_Y X$.

\begin{proposition}[5.4.2]
\label{I.5.4.2-fr}
Soit $S\to T$ un morphisme séparé. Si $X$ et $Y$ sont deux $S$-préschémas,
l'immersion canonique $X\times_S Y\to X\times_T Y$
(\hyperref[I.5.3.10-fr]{5.3.10}) est fermée.
\end{proposition}

En effet, si on se reporte au diagramme
(\hyperref[I.5.3.5.1-fr]{5.3.5.1}), on voit que $(p,q)_T$ peut être considéré
comme obtenu à partir de $\Delta_{S|T}$ par l'extension
$f\times_T g:X\times_T Y\to S\times_T S$ du préschéma de base $S\times_T S$~;
la proposition résulte alors de (\hyperref[I.4.3.2-fr]{4.3.2}).

\begin{corollary}[5.4.3]
\label{I.5.4.3-fr}
Soient $Y$ un $S$-schéma, $f:X\to Y$ un $S$-morphisme. Alors le morphisme
graphe $\Gamma_f:X\to X\times_S Y$ (\hyperref[I.5.3.11-fr]{5.3.11}) est une
immersion fermée.
\end{corollary}

C'est le cas particulier de (\hyperref[I.5.4.2-fr]{5.4.2}) où l'on remplace
$S$ par $Y$ et $T$ par $S$.

\begin{corollary}[5.4.4]
\label{I.5.4.4-fr}
Soient $f:X\to Y$, $g:Y\to Z$ deux morphismes, $g$ étant séparé. Si $g\circ f$
est une immersion fermée, il en est de même de $f$.
\end{corollary}

La démonstration à partir de (\hyperref[I.5.4.3-fr]{5.4.3}) est la même que
celle de (\hyperref[I.5.3.13-fr]{5.3.13}) à partir de
(\hyperref[I.5.3.11-fr]{5.3.11}).

\oldpage[I]{136}
\begin{corollary}[5.4.5]
\label{I.5.4.5-fr}
Soient $Z$ un $S$-schéma, $j:X\to Y$, $g:X\to Z$ deux $S$-morphismes. Si $j$
est une immersion fermée, il en est de même de
$(j,g)_S:X\to Y\times_S Z$.
\end{corollary}

La démonstration à partir de (\hyperref[I.5.4.4-fr]{5.4.4}) est la même que
celle de (\hyperref[I.5.3.14-fr]{5.3.14}) à partir de
(\hyperref[I.5.3.13-fr]{5.3.13}).

\begin{corollary}[5.4.6]
\label{I.5.4.6-fr}
Si $X$ est un $S$-schéma, toute $S$-section de $X$
(\hyperref[I.2.5.5-fr]{2.5.5}) est une immersion fermée.
\end{corollary}

Si $\varphi:X\to S$ est le morphisme structural, $\psi:S\to X$ une
$S$-section de $X$, il suffit d'appliquer
(\hyperref[I.5.4.5-fr]{5.4.5}) à $\varphi\circ\psi=1_S$.

\begin{corollary}[5.4.7]
\label{I.5.4.7-fr}
Soient $S$ un préschéma intègre, $s$ son point générique, $X$ un $S$-schéma.
Si deux $S$-sections $f$, $g$ de $X$ sont telles que $f(s)=g(s)$, alors
$f=g$.
\end{corollary}

En effet, si $x=f(s)=g(s)$, les homomorphismes $k(x)\to k(s)$ correspondant à
$f$ et $g$ sont nécessairement identiques. Si $h=(f,g)_S$, on en déduit
(\hyperref[I.5.3.17-fr]{5.3.17}) que $h(s)$ appartient à la diagonale
$Z=\Delta_X(X)$~; mais comme $S=\overline{\{s\}}$ et que $Z$ est fermée par
hypothèse, on a $h(S)\subset Z$. Il résulte alors de
(\hyperref[I.5.2.2-fr]{5.2.2}) que $h$ se factorise en
$S\to Z\to X\times_S X$, et on en conclut que $f=g$ par définition de la
diagonale.

\begin{remark}[5.4.8]
\label{I.5.4.8-fr}
Si l'on suppose réciproquement que la conclusion de
(\hyperref[I.5.4.3-fr]{5.4.3}) est vérifiée lorsque $f=1_Y$, on en conclut
que $Y$ est séparé au-dessus de $S$~; de même, si on suppose que la conclusion
de (\hyperref[I.5.4.5-fr]{5.4.5}) s'applique aux deux morphismes
\[
  Y\xrightarrow{\Delta_Y}Y\times_Z Y\xrightarrow{p_1}Y,
\]
on en tire que $\Delta_Y$ est une immersion fermée, donc que $Y$ est séparé
au-dessus de $Z$~; enfin, la validité de la conclusion de
(\hyperref[I.5.4.6-fr]{5.4.6}) pour la $Y$-section $\Delta_Y$ du
$Y$-préschéma $Y\times_S Y\to Y$, implique que $Y$ est séparé au-dessus de
$S$.
\end{remark}

\subsection{Critères de séparation.}
\label{subsection:I.5.5-fr}

\begin{proposition}[5.5.1]
\label{I.5.5.1-fr}
\begin{enumerate}
  \item[\rm(i)] Tout monomorphisme de préschémas (et en particulier toute
    immersion) est un morphisme séparé.
  \item[\rm(ii)] Le composé de deux morphismes séparés est séparé.
  \item[\rm(iii)] Si $f:X\to X'$, $g:Y\to Y'$ sont deux $S$-morphismes
    séparés, $f\times_S g$ est séparé.
  \item[\rm(iv)] Si $f:X\to Y$ est un $S$-morphisme séparé, le
    $S'$-morphisme $f_{(S')}$ est séparé pour toute extension $S'\to S$ du
    préschéma de base.
  \item[\rm(v)] Si le composé $g\circ f$ de deux morphismes est séparé, $f$
    est séparé.
  \item[\rm(vi)] Pour qu'un morphisme $f$ soit séparé, il faut et il suffit
    que $f_{\mathrm{red}}$ (\hyperref[I.5.1.5-fr]{5.1.5}) le soit.
\end{enumerate}
\end{proposition}

(i) résulte aussitôt de (\hyperref[I.5.3.8-fr]{5.3.8}). Si $f:X\to Y$,
$g:Y\to Z$ sont deux morphismes, le diagramme
\[
  \xymatrix{
    X\ar[rr]^{\Delta_{X|Z}}\ar[dr]_{\Delta_{X|Y}} & &
    X\times_Z X\\
    & X\times_Y X\ar[ur]_j
  }
  \tag{5.5.1.1}\label{I.5.5.1.1-fr}
\]
où $j$ désigne l'immersion canonique (\hyperref[I.5.3.10-fr]{5.3.10}) est
commutatif, comme on le vérifie aussitôt. Si $f$ et $g$ sont séparés,
$\Delta_{X|Y}$ est une immersion fermée par définition, et $j$ est une
immersion fermée en vertu de (\hyperref[I.5.4.2-fr]{5.4.2}), donc
$\Delta_{X|Z}$ est une immersion fermée par
(\hyperref[I.4.2.4-fr]{4.2.4}), ce qui

\oldpage[I]{137}
prouve (ii). Vu (i) et (ii), (iii) et (iv) sont équivalentes
(\hyperref[I.3.5.1-fr]{3.5.1}), et il suffit de démontrer (iv). Or,
$X_{(S')}\times_{Y_{(S')}}X_{(S')}$ s'identifie canoniquement à
$(X\times_Y X)\times_Y Y_{(S')}$ en vertu de
(\hyperref[I.3.3.11-fr]{3.3.11}) et de
(\hyperref[I.3.3.9.1-fr]{3.3.9.1}), et on vérifie aussitôt que le morphisme
diagonal $\Delta_{X_{(S')}}$ s'identifie alors à
$\Delta_X\times_Y 1_{Y_{(S')}}$~; la proposition résulte donc de
(\hyperref[I.4.3.1-fr]{4.3.1}).

Pour établir (v), considérons, comme dans
(\hyperref[I.5.3.13-fr]{5.3.13}) la factorisation
$X\xrightarrow{\Gamma_f}X\times_Z Y\xrightarrow{p_2}Y$ de $f$, en remarquant
que $p_2=(g\circ f)\times_Z 1_Y$~; l'hypothèse que $g\circ f$ est séparé
entraîne que $p_2$ est séparé par (iii) et (i), et comme $\Gamma_f$ est une
immersion, $\Gamma_f$ est séparé par (i), donc $f$ est séparé par (ii).
Enfin, pour démontrer (vi), rappelons que les préschémas
$X_{\mathrm{red}}\times_{Y_{\mathrm{red}}}X_{\mathrm{red}}$ et
$X_{\mathrm{red}}\times_Y X_{\mathrm{red}}$ s'identifient canoniquement
(\hyperref[I.5.1.7-fr]{5.1.7})~; si on désigne par $j$ l'injection
$X_{\mathrm{red}}\to X$, le diagramme
\[
  \xymatrix{
    X_{\mathrm{red}}\ar[r]^{\Delta_{X_{\mathrm{red}}}}\ar[d]_j &
    X_{\mathrm{red}}\times_Y X_{\mathrm{red}}\ar[d]^{j\times_Y j}\\
    X\ar[r]_{\Delta_X} & X\times_Y X
  }
\]
est commutatif (\hyperref[I.5.3.15-fr]{5.3.15}), et la proposition résulte de
ce que les flèches verticales sont des homéomorphismes des espaces
sous-jacents (\hyperref[I.4.3.1-fr]{4.3.1}).

\begin{corollary}[5.5.2]
\label{I.5.5.2-fr}
Si $f:X\to Y$ est séparé, la restriction de $f$ à tout sous-préschéma de $X$
est séparée.
\end{corollary}

Cela résulte de (\hyperref[I.5.5.1-fr]{5.5.1, (i) et (ii)}).

\begin{corollary}[5.5.3]
\label{I.5.5.3-fr}
Si $X$, $Y$ sont deux $S$-préschémas tels que $Y$ soit séparé au-dessus de
$S$, $X\times_S Y$ est séparé au-dessus de $X$.
\end{corollary}

C'est un cas particulier de (\hyperref[I.5.5.1-fr]{5.5.1, (iv)}).

\begin{proposition}[5.5.4]
\label{I.5.5.4-fr}
Soit $X$ un préschéma, et supposons que son espace sous-jacent soit réunion
d'une famille finie de parties fermées $X_k$ $(1\leq k\leq n)$~; on considère
pour chaque $k$ le sous-préschéma réduit de $X$ ayant $X_k$ pour espace
sous-jacent (\hyperref[I.5.2.1-fr]{5.2.1}) et on le note encore $X_k$. Soit
$f:X\to Y$ un morphisme, et pour chaque $k$, soit $Y_k$ une partie fermée de
$Y$ telle que $f(X_k)\subset Y_k$~; on note encore $Y_k$ le sous-préschéma
réduit de $Y$ ayant $Y_k$ pour espace sous-jacent, de sorte que la restriction
$X_k\to Y$ de $f$ à $X_k$ se factorise en
$X_k\xrightarrow{f_k}Y_k\to Y$ (\hyperref[I.5.2.2-fr]{5.2.2}). Pour que $f$
soit séparé, il faut et il suffit que les $f_k$ le soient.
\end{proposition}

La nécessité résulte de
(\hyperref[I.5.5.1-fr]{5.5.1, (i), (ii) et (v)}). Inversement, si la condition
de l'énoncé est satisfaite, chacune des restrictions $X_k\to Y$ de $f$ est
séparée (\hyperref[I.5.5.1-fr]{5.5.1, (i) et (ii)})~; si $p_1$, $p_2$ sont les
projections de $X\times_Y X$, le sous-espace $\Delta_{X_k}(X_k)$ s'identifie
au sous-espace $\Delta_X(X)\cap p_1^{-1}(X_k)$ de l'espace sous-jacent à
$X\times_Y X$ (\hyperref[I.5.3.16-fr]{5.3.16})~; ces sous-espaces étant
fermés dans $X\times_Y X$, il en est de même de leur réunion $\Delta_X(X)$.

Supposons en particulier que les $X_k$ soient les
\emph{composantes irréductibles} de $X$~; alors on peut supposer que les $Y_k$
sont des composantes irréductibles de $Y$
(\hyperref[0.2.1.5-fr]{0, 2.1.5})~; la prop.
(\hyperref[I.5.5.4-fr]{5.5.4}) ramène donc dans ce cas la notion de séparation
au cas des préschémas \emph{intègres} (\hyperref[I.2.1.7-fr]{2.1.7}).

\oldpage[I]{138}

\begin{proposition}[5.5.5]
\label{I.5.5.5-fr}
Soit $(Y_\lambda)$ un recouvrement ouvert d'un préschéma $Y$~; pour qu'un
morphisme $f:X\to Y$ soit séparé, il faut et il suffit que chacune de ses
restrictions $f^{-1}(Y_\lambda)\to Y_\lambda$ soit séparée.
\end{proposition}

Si on pose $X_\lambda=f^{-1}(Y_\lambda)$, tout revient, compte tenu de
(\hyperref[I.4.2.4-fr]{4.2.4, b)}) et de l'identité des produits
$X_\lambda\times_Y X_\lambda$ et
$X_\lambda\times_{Y_\lambda}X_\lambda$
(\hyperref[I.3.2.5-fr]{3.2.5}), à prouver que les
$X_\lambda\times_Y X_\lambda$ forment un recouvrement de $X\times_Y X$. Or,
si l'on pose $Y_{\lambda\mu}=Y_\lambda\cap Y_\mu$ et
$X_{\lambda\mu}=X_\lambda\cap X_\mu=f^{-1}(Y_{\lambda\mu})$,
$X_\lambda\times_Y X_\mu$ s'identifie au produit
$X_{\lambda\mu}\times_{Y_{\lambda\mu}}X_{\lambda\mu}$
(\hyperref[I.3.2.6.4-fr]{3.2.6.4}), donc aussi à
$X_{\lambda\mu}\times_Y X_{\lambda\mu}$
(\hyperref[I.3.2.5-fr]{3.2.5}), et finalement à un ouvert de
$X_\lambda\times_Y X_\lambda$, ce qui établit notre assertion
(\hyperref[I.3.2.7-fr]{3.2.7}).

La prop. (\hyperref[I.5.5.4-fr]{5.5.4}) permet, en prenant un recouvrement de
$Y$ par des ouverts affines, de ramener l'étude des morphismes séparés à celle
des morphismes séparés à valeurs dans des schémas affines.

\begin{proposition}[5.5.6]
\label{I.5.5.6-fr}
Soient $Y$ un schéma affine, $X$ un préschéma, $(U_\alpha)$ un recouvrement de
$X$ par des ouverts affines. Pour qu'un morphisme $f:X\to Y$ soit séparé, il
faut et il suffit que, pour tout couple d'indices $(\alpha,\beta)$,
$U_\alpha\cap U_\beta$ soit un ouvert affine, et que l'anneau
$\Gamma(U_\alpha\cap U_\beta,\mathscr{O}_X)$ soit engendré par la réunion des
images canoniques des anneaux $\Gamma(U_\alpha,\mathscr{O}_X)$ et
$\Gamma(U_\beta,\mathscr{O}_X)$.
\end{proposition}

Les $U_\alpha\times_Y U_\beta$ forment un recouvrement ouvert de
$X\times_Y X$ (\hyperref[I.3.2.7-fr]{3.2.7})~; en désignant par $p$ et $q$ les
projections de $X\times_Y X$, on a
\[
  \begin{aligned}
    \Delta_X^{-1}(U_\alpha\times_Y U_\beta)
      &=\Delta_X^{-1}\bigl(p^{-1}(U_\alpha)\cap q^{-1}(U_\beta)\bigr)\\
      &=\Delta_X^{-1}\bigl(p^{-1}(U_\alpha)\bigr)
        \cap\Delta_X^{-1}\bigl(q^{-1}(U_\beta)\bigr)
        =U_\alpha\cap U_\beta~;
  \end{aligned}
\]
tout revient donc à exprimer que la restriction de $\Delta_X$ à
$U_\alpha\cap U_\beta$ est une immersion fermée dans
$U_\alpha\times_Y U_\beta$. Or, cette restriction n'est autre que
$(j_\alpha,j_\beta)_Y$, en désignant par $j_\alpha$ (resp. $j_\beta$) le
morphisme d'injection de $U_\alpha\cap U_\beta$ dans $U_\alpha$ (resp.
$U_\beta$), ainsi qu'il résulte des définitions. Comme
$U_\alpha\times_Y U_\beta$ est un schéma affine dont l'anneau est
canoniquement isomorphe à
$\Gamma(U_\alpha,\mathscr{O}_X)
 \otimes_{\Gamma(Y,\mathscr{O}_Y)}
 \Gamma(U_\beta,\mathscr{O}_X)$
(\hyperref[I.3.2.2-fr]{3.2.2}), on voit que $U_\alpha\cap U_\beta$ doit être
un schéma affine et que l'application
$h_\alpha\otimes h_\beta\to h_\alpha h_\beta$ de l'anneau
$A(U_\alpha\times_Y U_\beta)$ dans
$\Gamma(U_\alpha\cap U_\beta,\mathscr{O}_X)$ doit être surjective
(\hyperref[I.4.2.3-fr]{4.2.3}), ce qui achève la démonstration.

\begin{corollary}[5.5.7]
\label{I.5.5.7-fr}
Un schéma affine est séparé (et est par suite un schéma, ce qui justifie la
terminologie de (\hyperref[I.5.4.1-fr]{5.4.1})).
\end{corollary}

\begin{corollary}[5.5.8]
\label{I.5.5.8-fr}
Soit $Y$ un schéma affine~; pour que $f:X\to Y$ soit un morphisme séparé, il
faut et il suffit que $X$ soit séparé (autrement dit, que $X$ soit un schéma).
\end{corollary}

On observera en effet que le critère de
(\hyperref[I.5.5.6-fr]{5.5.6}) ne dépend pas de $f$.

\begin{corollary}[5.5.9]
\label{I.5.5.9-fr}
Pour qu'un morphisme $f:X\to Y$ soit séparé, il faut que pour tout ouvert $U$
sur lequel $Y$ induit un préschéma séparé, le préschéma induit $f^{-1}(U)$
soit séparé, et il suffit qu'il en soit ainsi pour tout ouvert affine
$U\subset Y$.
\end{corollary}

La nécessité de la condition résulte de
(\hyperref[I.5.5.4-fr]{5.5.4}) et
(\hyperref[I.5.5.1-fr]{5.5.1, (ii)})~; la suffisance résulte de
(\hyperref[I.5.5.4-fr]{5.5.4}) et
(\hyperref[I.5.5.8-fr]{5.5.8}), compte tenu de l'existence de recouvrements
ouverts affines de $Y$.

En particulier, si $X$ et $Y$ sont des schémas affines, \emph{tout} morphisme
$X\to Y$ est séparé.

\begin{proposition}[5.5.10]
\label{I.5.5.10-fr}
Soient $Y$ un schéma, $f:X\to Y$ un morphisme. Pour tout ouvert affine $U$ de
$X$ et tout ouvert affine $V$ de $Y$, $U\cap f^{-1}(V)$ est affine.
\end{proposition}

Soient $p_1$, $p_2$ les projections de $X\times_{\mathbf{Z}}Y$~; le
sous-espace $U\cap f^{-1}(V)$ est l'image par $p_1$ de
$\Gamma_f(X)\cap p_1^{-1}(U)\cap p_2^{-1}(V)$. Or,
$p_1^{-1}(U)\cap p_2^{-1}(V)$ s'identifie à l'espace sous-jacent au

\oldpage[I]{139}
préschéma $U\times_{\mathbf{Z}}V$
(\hyperref[I.3.2.7-fr]{3.2.7}), et est par suite un schéma affine
(\hyperref[I.3.2.2-fr]{3.2.2})~; comme $\Gamma_f(X)$ est fermé dans
$X\times_{\mathbf{Z}}Y$ (\hyperref[I.5.4.3-fr]{5.4.3}),
$\Gamma_f(X)\cap p_1^{-1}(U)\cap p_2^{-1}(V)$ est fermé dans
$U\times_{\mathbf{Z}}V$, et par suite le préschéma induit par le
sous-préschéma de $X\times_{\mathbf{Z}}Y$ associé à $\Gamma_f$
(\hyperref[I.4.2.1-fr]{4.2.1}), sur la partie ouverte
$\Gamma_f(X)\cap p_1^{-1}(U)\cap p_2^{-1}(V)$ de son espace sous-jacent, est
un sous-préschéma fermé d'un schéma affine, donc un schéma affine
(\hyperref[I.4.2.3-fr]{4.2.3}). La proposition résulte alors du fait que
$\Gamma_f$ est une immersion.

\begin{examples}[5.5.11]
\label{I.5.5.11-fr}
Le préschéma de l'exemple (\hyperref[I.2.3.2-fr]{2.3.2})
(<<~droite projective sur un corps $K$~>>) est \emph{séparé}, car pour le
recouvrement $(X_1,X_2)$ de $X$ par des ouverts affines,
$X_1\cap X_2=U_{12}$ est affine et $\Gamma(U_{12},\mathscr{O}_X)$, anneau des
fractions rationnelles de la forme $f(s)/s^m$ avec $f\in K[s]$, est engendré
par $K[s]$ et par $1/s$, donc les conditions de
(\hyperref[I.5.5.6-fr]{5.5.6}) sont vérifiées.

Avec le même choix de $X_1$, $X_2$, $U_{12}$ et $U_{21}$ que dans l'exemple
(\hyperref[I.2.3.2-fr]{2.3.2}), prenons cette fois pour $u_{12}$
l'isomorphisme qui à $f(s)$ fait correspondre $f(t)$~; on obtient cette fois
par recollement un préschéma \emph{intègre non séparé} $X$, car la première
condition de (\hyperref[I.5.5.6-fr]{5.5.6}) est vérifiée, mais non la seconde.
Il est immédiat ici que
$\Gamma(X,\mathscr{O}_X)\to\Gamma(X_1,\mathscr{O}_X)=K[s]$ est un
isomorphisme~; l'isomorphisme réciproque définit un morphisme
$f:X\to\operatorname{Spec}(K[s])$ qui est surjectif, et pour tout
$y\in\operatorname{Spec}(K[s])$ tel que $\mathfrak{j}_y\neq(0)$,
$f^{-1}(y)$ est réduit à un point, mais pour $\mathfrak{j}_y=(0)$,
$f^{-1}(y)$ se compose de \emph{deux} points distincts (on dit que $X$ est la
<<~droite affine sur $K$, où le point $0$ est dédoublé~>>).

On peut aussi donner des exemples où \emph{aucune} des deux conditions de
(\hyperref[I.5.5.6-fr]{5.5.6}) n'est vérifiée. Remarquons d'abord que dans le
spectre premier $Y$ de l'anneau de polynômes $A=K[s,t]$ à deux indéterminées
sur un corps $K$, l'ouvert $U$ réunion de $D(s)$ et de $D(t)$ \emph{n'est pas
un ouvert affine}. En effet, si $z$ est une section de $\mathscr{O}_Y$
au-dessus de $U$, il existe deux entiers $m\geq0$, $n\geq0$ tels que $s^m z$
et $t^n z$ soient les restrictions à $U$ de polynômes en $s$ et $t$
(\hyperref[I.1.4.1-fr]{1.4.1}), ce qui n'est évidemment possible que si la
section $z$ se prolonge en une section au-dessus de $Y$ tout entier,
identifiée à un polynôme en $s$ et $t$. Si $U$ était un ouvert affine, le
morphisme d'injection $U\to Y$ serait donc un isomorphisme
(\hyperref[I.1.7.3-fr]{1.7.3}), ce qui est absurde puisque $U\neq Y$.

Cela étant, prenons deux schémas affines $Y_1$, $Y_2$, spectres premiers des
anneaux $A_1=K[s_1,t_1]$, $A_2=K[s_2,t_2]$~; prenons
$U_{12}=D(s_1)\cup D(t_1)$, $U_{21}=D(s_2)\cup D(t_2)$, et pour $u_{12}$ la
restriction à $U_{21}$ de l'isomorphisme $Y_2\to Y_1$ correspondant à
l'isomorphisme d'anneaux qui à $f(s_1,t_1)$ fait correspondre $f(s_2,t_2)$~;
on a ainsi un exemple où aucune des conditions de
(\hyperref[I.5.5.6-fr]{5.5.6}) n'est satisfaite (le préschéma intègre ainsi
obtenu est dit <<~plan affine sur $K$, où le point $0$ est dédoublé~>>).
\end{examples}

\begin{remark}[5.5.12]
\label{I.5.5.12-fr}
\normalfont Étant donnée une propriété \textbf{P} de morphismes de préschémas,
considérons les propositions suivantes~:
\begin{enumerate}
  \item[\rm(i)] \emph{Toute immersion fermée possède la propriété \textbf{P}.}
  \item[\rm(ii)] \emph{Le composé de deux morphismes possédant la propriété
    \textbf{P} possède la propriété \textbf{P}.}
  \item[\rm(iii)] \emph{Si $f:X\to X'$, $g:Y\to Y'$ sont deux $S$-morphismes
    possédant la propriété \textbf{P}, $f\times_S g$ possède la propriété
    \textbf{P}.}

\oldpage[I]{140}
  \item[\rm(iv)] \emph{Si $f:X\to Y$ est un $S$-morphisme possédant la
    propriété \textbf{P}, tout $S'$-morphisme $f_{(S')}$ obtenu par une
    extension $S'\to S$ du préschéma de base, possède la propriété
    \textbf{P}.}
  \item[\rm(v)] \emph{Si le composé $g\circ f$ de deux morphismes
    $f:X\to Y$, $g:Y\to Z$ possède la propriété \textbf{P}, et si $g$ est
    séparé, $f$ possède la propriété \textbf{P}.}
  \item[\rm(vi)] \emph{Si un morphisme $f:X\to Y$ possède la propriété
    \textbf{P}, il en est de même de $f_{\mathrm{red}}$
    (\hyperref[I.5.1.5-fr]{5.1.5}).}
\end{enumerate}

Dans ces conditions, si on suppose (i) et (ii) vérifiées, (iii) et (iv) sont
\emph{équivalentes}, et (v) et (vi) sont des \emph{conséquences} de (i), (ii)
et (iii).

La première assertion a déjà été démontrée
(\hyperref[I.3.5.1-fr]{3.5.1}). Considérons la factorisation
(\hyperref[I.5.3.13-fr]{5.3.13}) de $f$ en
$X\xrightarrow{\Gamma_f}X\times_Z Y\xrightarrow{p_2}Y$~; la relation
$p_2=(g\circ f)\times_Z 1_Y$ montre que si $g\circ f$ possède la propriété
\textbf{P}, il en est de même de $p_2$ en vertu de (iii)~; si $g$ est séparé,
$\Gamma_f$ est une immersion fermée
(\hyperref[I.5.4.3-fr]{5.4.3}), donc possède aussi la propriété \textbf{P} par
(i)~; enfin, en vertu de (ii), $f$ possède la propriété \textbf{P}.

Enfin, considérons le diagramme commutatif
\[
  \xymatrix{
    X_{\mathrm{red}}\ar[r]^{f_{\mathrm{red}}}\ar[d] &
    Y_{\mathrm{red}}\ar[d]\\
    X\ar[r]^f &
    Y
  }
\]
où les flèches verticales sont des immersions fermées
(\hyperref[I.5.1.5-fr]{5.1.5}), donc possèdent la propriété \textbf{P} par
(i). L'hypothèse que $f$ possède la propriété \textbf{P} entraîne donc par
(ii) que
$X_{\mathrm{red}}\xrightarrow{f_{\mathrm{red}}}Y_{\mathrm{red}}\to Y$
possède la propriété \textbf{P}~; enfin, comme une immersion fermée est
séparée (\hyperref[I.5.5.1-fr]{5.5.1, (i)}), $f_{\mathrm{red}}$ possède la
propriété \textbf{P} en vertu de (v).

On notera que si on considère les propositions~:
\begin{enumerate}
  \item[\rm(i')] \emph{Toute immersion possède la propriété \textbf{P}.}
  \item[\rm(v')] \emph{Si $g\circ f$ possède la propriété \textbf{P}, il en
    est de même de $f$~;}
\end{enumerate}
alors le raisonnement fait ci-dessus montre que (v') est conséquence de (i'),
(ii) et (iii).
\end{remark}

\begin{env}[5.5.13]
\label{I.5.5.13-fr}
On notera que (v) et (vi) sont encore conséquences de (i), (iii) et
\begin{enumerate}
  \item[\rm(ii')] \emph{Si $j:X\to Y$ est une immersion fermée et
    $g:Y\to Z$ un morphisme possédant la propriété \textbf{P}, alors
    $g\circ j$ possède la propriété \textbf{P}.}
\end{enumerate}
De même, (v') est conséquence de (i'), (iii) et
\begin{enumerate}
  \item[\rm(ii'')] \emph{Si $j:X\to Y$ est une immersion et $g:Y\to Z$ un
    morphisme possédant la propriété \textbf{P}, alors $g\circ j$ possède la
    propriété \textbf{P}.}
\end{enumerate}
Cela résulte en effet aussitôt des raisonnements de
(\hyperref[I.5.5.12-fr]{5.5.12}).
\end{env}
